% GENERATED FILE. Edit canonical lessons, then run npm run generate:books.
% canonical-source-hash: fnv1a64:99561fd7bcd6e1dd
% canonical-lessons: SA-C295-cetana, SA-C295-utthapayati, SA-C295-sibika, SA-C295-apatkalah, SA-C295-avalambate

\chapter{Help for the Injured}
\label{ch:sa-help-for-the-injured}
\hlchaptermodality{295}

\begin{chapteropening}
\textbf{By the end of this chapter:} \emph{I can say how people helped someone who was hurt.}

Complete the last lesson of chapter 295: I can say how people helped someone who was hurt.
\end{chapteropening}

\section[\texorpdfstring{\sk{चेतना} (cetanā)}{चेतना (cetanā)}]{\sk{चेतना} (cetanā) — consciousness}

\label{lesson:SA-C295-cetana}



\begin{quote}
Before the new one: say the Sanskrit for injured, then the Sanskrit for faints.
\end{quote}

\subsection*{You'll want to know: \sk{चेतना}}
\textbf{\sk{चेतना}} — \emph{cetanā} — \textquotedblleft{}consciousness\textquotedblright{}.

\begin{quote}
\sk{आहतस्य} \sk{चेतना} \sk{शनैः} \sk{आगच्छति।}
\end{quote}

\begin{quote}
\emph{āhatasya cetanā śanaiḥ āgacchati.}
\end{quote}

\begin{quote}
The injured man slowly comes round.
\end{quote}

A feminine noun.

\begin{cousinweb}
\textbf{\sk{चेतना}} (\emph{cetanā}) comes from the root √cit, \textquotedblleft{}to perceive, attend to\textquotedblright{}; MW gives \emph{cetanā} as consciousness and understanding (MW).
\end{cousinweb}

\begin{grammarlens}[title={What you've built}]
One more word for telling what happened.
\end{grammarlens}

\subsection*{Your turn}
\begin{itemize}
\raggedright
  \item \emph{Say it:} \emph{cetanā}
  \item \emph{Say it:} \emph{cetanā}, once more
  \item \emph{Say it:} \emph{āhatasya cetanā śanaiḥ āgacchati.} — The injured man slowly comes round.
  \item \emph{Recall:} say the Sanskrit for injured, then the Sanskrit for faints, then say \emph{cetanā} again
\end{itemize}

\begin{tcolorbox}[breakable,colback=teal!4,colframe=teal!35!black,title={Before you move on}]
What does \sk{चेतना} mean? (\textquotedblleft{}Consciousness\textquotedblright{}.) Say it once more.
\end{tcolorbox}
\section[\texorpdfstring{\sk{उत्थापयति} (utthāpayati)}{उत्थापयति (utthāpayati)}]{\sk{उत्थापयति} (utthāpayati) — lifts up; raises}

\label{lesson:SA-C295-utthapayati}



\begin{quote}
Before the new one: say the Sanskrit for faints, then the Sanskrit for consciousness.
\end{quote}

\subsection*{You'll want to know: \sk{उत्थापयति}}
\textbf{\sk{उत्थापयति}} — \emph{utthāpayati} — \textquotedblleft{}lifts up; raises\textquotedblright{}.

\begin{quote}
\sk{जनाः} \sk{आहतम्} \sk{उत्थापयन्ति।}
\end{quote}

\begin{quote}
\emph{janāḥ āhatam utthāpayanti.}
\end{quote}

\begin{quote}
People lift up the injured man.
\end{quote}

Past (he …): \sk{सः} \sk{आहतम्} \sk{उत्थापितवान्}, \emph{saḥ āhatam utthāpitavān}.

\begin{cousinweb}
\textbf{\sk{उत्थापयति}} (\emph{utthāpayati}) is built from \emph{ud}, \textquotedblleft{}up\textquotedblright{}, and the causative of the root √sthā, \textquotedblleft{}to stand\textquotedblright{}: to make stand up, to raise (MW).
\end{cousinweb}

\begin{grammarlens}[title={What you've built}]
One more word for telling what happened.
\end{grammarlens}

\subsection*{Your turn}
\begin{itemize}
\raggedright
  \item \emph{Say it:} \emph{utthāpayati}
  \item \emph{Say it:} \emph{utthāpayati}, once more
  \item \emph{Say it:} \emph{janāḥ āhatam utthāpayanti.} — People lift up the injured man.
  \item \emph{Recall:} say the Sanskrit for faints, then the Sanskrit for consciousness, then say \emph{utthāpayati} again
\end{itemize}

\begin{tcolorbox}[breakable,colback=teal!4,colframe=teal!35!black,title={Before you move on}]
What does \sk{उत्थापयति} mean? (\textquotedblleft{}Lifts up; raises\textquotedblright{}.) Say it once more.
\end{tcolorbox}
\section[\texorpdfstring{\sk{शिबिका} (śibikā)}{शिबिका (śibikā)}]{\sk{शिबिका} (śibikā) — a litter; a stretcher}

\label{lesson:SA-C295-sibika}



\begin{quote}
Before the new one: say the Sanskrit for consciousness, then the Sanskrit for lifts up.
\end{quote}

\subsection*{You'll want to know: \sk{शिबिका}}
\textbf{\sk{शिबिका}} — \emph{śibikā} — \textquotedblleft{}a litter; a stretcher\textquotedblright{}.

\begin{quote}
\sk{ते} \sk{आहतं} \sk{शिबिकया} \sk{नयन्ति।}
\end{quote}

\begin{quote}
\emph{te āhataṃ śibikayā nayanti.}
\end{quote}

\begin{quote}
They carry the injured man on a litter.
\end{quote}

A feminine noun.

\begin{cousinweb}
\textbf{\sk{शिबिका}} (\emph{śibikā}) MW gives \emph{śibikā} as a palanquin, a litter or a bier, without a derivation, so its origin is uncertain; a stretcher is its modern use (MW).
\end{cousinweb}

\begin{grammarlens}[title={What you've built}]
One more word for telling what happened.
\end{grammarlens}

\subsection*{Your turn}
\begin{itemize}
\raggedright
  \item \emph{Say it:} \emph{śibikā}
  \item \emph{Say it:} \emph{śibikā}, once more
  \item \emph{Say it:} \emph{te āhataṃ śibikayā nayanti.} — They carry the injured man on a litter.
  \item \emph{Recall:} say the Sanskrit for consciousness, then the Sanskrit for lifts up, then say \emph{śibikā} again
\end{itemize}

\begin{tcolorbox}[breakable,colback=teal!4,colframe=teal!35!black,title={Before you move on}]
What does \sk{शिबिका} mean? (\textquotedblleft{}A litter; a stretcher\textquotedblright{}.) Say it once more.
\end{tcolorbox}
\section[\texorpdfstring{\sk{आपत्कालः} (āpatkālaḥ)}{आपत्कालः (āpatkālaḥ)}]{\sk{आपत्कालः} (āpatkālaḥ) — an emergency}

\label{lesson:SA-C295-apatkalah}



\begin{quote}
Before the new one: say the Sanskrit for lifts up, then the Sanskrit for a litter.
\end{quote}

\subsection*{You'll want to know: \sk{आपत्कालः}}
\textbf{\sk{आपत्कालः}} — \emph{āpatkālaḥ} — \textquotedblleft{}an emergency\textquotedblright{}.

\begin{quote}
\sk{आपत्काले} \sk{वैद्यः} \sk{शीघ्रम्} \sk{आगच्छति।}
\end{quote}

\begin{quote}
\emph{āpatkāle vaidyaḥ śīghram āgacchati.}
\end{quote}

\begin{quote}
In an emergency the doctor comes quickly.
\end{quote}

A masculine noun.

\begin{cousinweb}
\textbf{\sk{आपत्कालः}} (\emph{āpatkālaḥ}) is built from \emph{āpad}, \textquotedblleft{}misfortune, distress\textquotedblright{}, and \emph{kāla}, \textquotedblleft{}time\textquotedblright{}; MW gives it as a time of distress (MW).
\end{cousinweb}

\begin{grammarlens}[title={What you've built}]
One more word for telling what happened.
\end{grammarlens}

\subsection*{Your turn}
\begin{itemize}
\raggedright
  \item \emph{Say it:} \emph{āpatkālaḥ}
  \item \emph{Say it:} \emph{āpatkālaḥ}, once more
  \item \emph{Say it:} \emph{āpatkāle vaidyaḥ śīghram āgacchati.} — In an emergency the doctor comes quickly.
  \item \emph{Recall:} say the Sanskrit for lifts up, then the Sanskrit for a litter, then say \emph{āpatkālaḥ} again
\end{itemize}

\begin{tcolorbox}[breakable,colback=teal!4,colframe=teal!35!black,title={Before you move on}]
What does \sk{आपत्कालः} mean? (\textquotedblleft{}An emergency\textquotedblright{}.) Say it once more.
\end{tcolorbox}
\section[\texorpdfstring{\sk{अवलम्बते} (avalambate)}{अवलम्बते (avalambate)}]{\sk{अवलम्बते} (avalambate) — leans on}

\label{lesson:SA-C295-avalambate}



\begin{quote}
Before the new one: say the Sanskrit for a litter, then the Sanskrit for an emergency.
\end{quote}

\subsection*{You'll want to know: \sk{अवलम्बते}}
\textbf{\sk{अवलम्बते}} — \emph{avalambate} — \textquotedblleft{}leans on\textquotedblright{}.

\begin{quote}
\sk{आहतः} \sk{मित्रस्य} \sk{स्कन्धम्} \sk{अवलम्बते।}
\end{quote}

\begin{quote}
\emph{āhataḥ mitrasya skandham avalambate.}
\end{quote}

\begin{quote}
The injured man leans on his friend's shoulder.
\end{quote}

Past (he …): \sk{सः} \sk{मित्रम्} \sk{अवलम्बितवान्}, \emph{saḥ mitram avalambitavān}.

\begin{cousinweb}
\textbf{\sk{अवलम्बते}} (\emph{avalambate}) is built from \emph{ava}, \textquotedblleft{}down\textquotedblright{}, and the root √lamb, \textquotedblleft{}to hang\textquotedblright{}; MW gives it as to hang down and to cling to or rest upon as a support (MW).
\end{cousinweb}

\begin{grammarlens}[title={What you've built}]
One more word for telling what happened.
\end{grammarlens}

\subsection*{Your turn}
\begin{itemize}
\raggedright
  \item \emph{Say it:} \emph{avalambate}
  \item \emph{Say it:} \emph{avalambate}, once more
  \item \emph{Say it:} \emph{āhataḥ mitrasya skandham avalambate.} — The injured man leans on his friend's shoulder.
  \item \emph{Recall:} say the Sanskrit for a litter, then the Sanskrit for an emergency, then say \emph{avalambate} again
\end{itemize}

\begin{tcolorbox}[breakable,colback=teal!4,colframe=teal!35!black,title={Before you move on}]
What does \sk{अवलम्बते} mean? (\textquotedblleft{}Leans on\textquotedblright{}.) Say it once more.
\end{tcolorbox}
